\documentclass[a4paper, 12pt]{article}
\usepackage[T2A]{fontenc}
\usepackage[utf8]{inputenc}
\usepackage[russian]{babel}
\usepackage{amsmath, amssymb}
\usepackage{physics}
\usepackage{bbm}
\usepackage{bm}
\usepackage{hyperref}
\usepackage{graphicx}
\usepackage[dvipsnames]{xcolor}
\usepackage{tikz}
\usetikzlibrary{calc,decorations.markings,patterns,positioning}
\usepackage{tikz-3dplot}
\usepackage{pgfplots}
\usepackage{subcaption}
\usepackage[verbose=true, a4paper, left=20mm, right=20mm, top=20mm,
bottom=20mm]{geometry}
\usepackage{chngcntr}
\usepackage{tocloft}
% Настройка математических операторов
\renewcommand{\tan}{\operatorname{tg}}
\renewcommand{\cot}{\operatorname{ctg}}
\renewcommand{\sinh}{\operatorname{sh}}
\renewcommand{\cosh}{\operatorname{ch}}
% Настройка нумерации уравнений по разделам (только для основной части)
\numberwithin{equation}{section}
% Настройка глубины оглавления
\setcounter{tocdepth}{3}
\setcounter{secnumdepth}{3}
% Настройка отступов для номеров разделов в оглавлении
\renewcommand{\cftsecnumwidth}{2em}
\renewcommand{\cftsubsecnumwidth}{3em}
\renewcommand{\cftsubsubsecnumwidth}{4em}
\title{Конспект лекций по курсу <<Квантовая Механика 1>>  \\ \large Лекция 8}
\author{Лектор: Грибанов Сергей Сергеевич}
\date{5 марта}
\begin{document}
\maketitle
\tableofcontents

\section{Уравнение Хартри}
\label{sec:hartree_equation}

Мы установили, что в приближении самосогласованного поля каждый электрон движется независимо в некотором эффективном центральном потенциале, создаваемом ядром и усреднённым распределением заряда остальных электронов. Перейдём теперь к построению конкретных уравнений, определяющих этот потенциал и соответствующие одночастичные состояния. Начнём с наиболее простого подхода, предложенного Хартри.

\subsection{Вариационный вывод}
\label{subsec:hartree_variational}

В методе Хартри волновая функция $Z$-электронного атома аппроксимируется простым произведением одночастичных волновых функций:
\begin{equation}
\label{eq:hartree_product}
\Psi(\bm{r}_1, \bm{r}_2, \dots, \bm{r}_Z) = \psi_{\nu_1}(\bm{r}_1) \psi_{\nu_2}(\bm{r}_2) \dots \psi_{\nu_Z}(\bm{r}_Z),
\end{equation}
где $\nu_i$ --- набор квантовых чисел, характеризующий состояние $i$-го электрона. Такая функция не антисимметрична относительно перестановки частиц. Принцип Паули в этом подходе учитывается не на уровне волновой функции, а на уровне выбора набора состояний: предполагается, что никакие два электрона не занимают одно и то же одночастичное состояние. Это приближение позволяет получить замкнутую систему уравнений для одночастичных функций и во многих случаях даёт результаты, близкие к результатам более точных методов.

Для нахождения наилучших функций $\psi_{\nu_i}$ используем вариационный принцип. Вычислим среднее значение гамильтониана \eqref{eq:many_electron_hamiltonian} на функции \eqref{eq:hartree_product}. Одночастичная часть даёт
\begin{equation}
\label{eq:hartree_one_electron}
\langle \hat{T} + \hat{V}_{\text{яд}} \rangle = \sum_{a=1}^{Z} \int \dd[3]{\bm{r}_a} \psi_{\nu_a}^*(\bm{r}_a) \left( -\frac{\hslash^2}{2m} \nabla_a^2 - \frac{Z e^2}{r_a} \right) \psi_{\nu_a}(\bm{r}_a).
\end{equation}
Энергия межэлектронного отталкивания вычисляется как
\begin{equation}
\label{eq:hartree_two_electron}
\langle \hat{V}_{\text{ээ}} \rangle = \frac{1}{2} \sum_{a \neq b} \iint \dd[3]{\bm{r}_a} \dd[3]{\bm{r}_b} \frac{e^2}{\abs{\bm{r}_a - \bm{r}_b}} \abs{\psi_{\nu_a}(\bm{r}_a)}^2 \abs{\psi_{\nu_b}(\bm{r}_b)}^2.
\end{equation}
Множитель $1/2$ компенсирует двойной учёт пар при суммировании по всем $a \neq b$. Полная средняя энергия равна сумме \eqref{eq:hartree_one_electron} и \eqref{eq:hartree_two_electron}:
\begin{equation}
\label{eq:hartree_energy}
\begin{split}
E[\psi] &= \sum_{a=1}^{Z} \int \dd[3]{\bm{r}_a} \psi_{\nu_a}^*(\bm{r}_a) \left( -\frac{\hslash^2}{2m} \nabla_a^2 - \frac{Z e^2}{r_a} \right) \psi_{\nu_a}(\bm{r}_a) \\
&\quad + \frac{1}{2} \sum_{a \neq b} \iint \dd[3]{\bm{r}_a} \dd[3]{\bm{r}_b} \frac{e^2}{\abs{\bm{r}_a - \bm{r}_b}} \abs{\psi_{\nu_a}(\bm{r}_a)}^2 \abs{\psi_{\nu_b}(\bm{r}_b)}^2.
\end{split}
\end{equation}
Варьируя этот функционал по $\psi_{\nu_a}^*(\bm{r}_a)$ с учётом условий нормировки
\begin{equation}
\int \dd[3]{\bm{r}_a} \abs{\psi_{\nu_a}(\bm{r}_a)}^2 = 1,
\end{equation}
и вводя соответствующие множители Лагранжа $\varepsilon_{\nu_a}$, получаем уравнения Хартри:
\begin{equation}
\label{eq:hartree_eq_general}
\left( -\frac{\hslash^2}{2m} \nabla_a^2 - \frac{Z e^2}{r_a} + \sum_{b \neq a} \int \dd[3]{\bm{r}_b} \frac{e^2 \abs{\psi_{\nu_b}(\bm{r}_b)}^2}{\abs{\bm{r}_a - \bm{r}_b}} \right) \psi_{\nu_a}(\bm{r}_a) = \varepsilon_{\nu_a} \psi_{\nu_a}(\bm{r}_a).
\end{equation}
Уравнения \eqref{eq:hartree_eq_general} представляют собой систему нелинейных интегро-дифференциальных уравнений: потенциал, действующий на $a$-й электрон, зависит от волновых функций всех остальных электронов.

\subsection{Эффективный потенциал и самосогласование}
\label{subsec:hartree_effective_potential}

Физическая интерпретация уравнения Хартри прозрачна. Каждый электрон движется в эффективном потенциале, который складывается из потенциала ядра и потенциала, создаваемого усреднённым распределением заряда всех остальных электронов:
\begin{equation}
\label{eq:hartree_potential}
V_a^{\text{H}}(\bm{r}_a) = - \frac{Z e^2}{r_a} + \sum_{b \neq a} \int \dd[3]{\bm{r}_b} \frac{e^2 \abs{\psi_{\nu_b}(\bm{r}_b)}^2}{\abs{\bm{r}_a - \bm{r}_b}}.
\end{equation}
Второе слагаемое удобно выразить через плотность заряда остальных электронов:
\begin{equation}
\label{eq:hartree_density}
\rho_a(\bm{r}) = e \sum_{b \neq a} \abs{\psi_{\nu_b}(\bm{r})}^2,
\end{equation}
так что
\begin{equation}
V_a^{\text{H}}(\bm{r}_a) = - \frac{Z e^2}{r_a} + \int \dd[3]{\bm{r}_b} \frac{\rho_a(\bm{r}_b)}{\abs{\bm{r}_a - \bm{r}_b}}.
\end{equation}
Этот потенциал не является центральным в общем случае, поскольку распределение заряда других электронов может не обладать сферической симметрией. Однако для атомов в основном состоянии можно ожидать, что усреднённое распределение заряда будет сферически симметричным. Действительно, если оболочки заполнены или усреднены по проекциям момента, плотность $\abs{\psi_{nlm}}^2$, просуммированная по всем $m$ от $-l$ до $l$, не зависит от углов:
\begin{equation}
\label{eq:spherical_density}
\sum_{m=-l}^{l} \abs{\psi_{nlm}(\bm{r})}^2 = \frac{2l+1}{4\pi} R_{nl}^2(r).
\end{equation}
Поэтому в приближении центрального поля потенциал $V_a^{\text{H}}(\bm{r}_a)$ становится функцией только $r_a$, что существенно упрощает задачу.

Требование самосогласованности состоит в том, что потенциал $V_a^{\text{H}}$, вычисленный по волновым функциям, должен совпадать с потенциалом, который использовался при решении уравнения Шрёдингера для этих функций. Поскольку потенциал зависит от функций, а функции --- от потенциала, задача решается методом последовательных приближений. В качестве нулевого приближения обычно берут волновые функции водородоподобного атома с зарядом ядра $Z$ или с некоторым эффективным зарядом $\tilde{Z}$. По этим функциям вычисляют распределение заряда и потенциал \eqref{eq:hartree_potential}. Затем решают уравнение Шрёдингера с этим потенциалом, находят новые волновые функции и повторяют процесс до тех пор, пока потенциал и функции не перестанут заметно изменяться от итерации к итерации. Такой итерационный процесс, если он сходится, приводит к самосогласованному решению.

Решения уравнения Хартри в центральном поле характеризуются квантовыми числами $n, l, m, \sigma$. Энергия $\varepsilon_{nl}$ зависит от $n$ и $l$, поскольку эффективный потенциал отличается от кулоновского. Порядок уровней определяется экранировкой: уровни с меньшим $l$ лежат ниже, так как соответствующие электроны сильнее притягиваются к ядру, глубже проникая в область слабой экранировки. Кратность вырождения уровня $\varepsilon_{nl}$ по проекциям момента и спина равна $2(2l+1)$.

\section{Метод Хартри---Фока}
\label{sec:hartree_fock_method}

Уравнение Хартри учитывает только прямое кулоновское взаимодействие электронов. Оно не включает обменные эффекты, связанные с антисимметрией полной волновой функции. Для более точного описания необходимо перейти к методу Хартри---Фока, который последовательно учитывает принцип Паули.

\subsection{Определитель Слэтера}
\label{subsec:slater_determinant}

Волновая функция совокупности фермионов должна быть антисимметричной относительно перестановки любой пары частиц. Для $Z$ электронов такой функцией является определитель Слэтера, построенный из одночастичных ортонормированных функций $\psi_{\nu_a}(\bm{r}_a, \sigma_a)$:
\begin{equation}
\label{eq:slater_det}
\Psi(1, 2, \dots, Z) = \frac{1}{\sqrt{Z!}}
\begin{vmatrix}
\psi_{\nu_1}(1) & \psi_{\nu_1}(2) & \dots & \psi_{\nu_1}(Z) \\
\psi_{\nu_2}(1) & \psi_{\nu_2}(2) & \dots & \psi_{\nu_2}(Z) \\
\vdots & \vdots & \ddots & \vdots \\
\psi_{\nu_Z}(1) & \psi_{\nu_Z}(2) & \dots & \psi_{\nu_Z}(Z)
\end{vmatrix},
\end{equation}
где для краткости введено обозначение $a \equiv (\bm{r}_a, \sigma_a)$. Определитель автоматически антисимметричен относительно перестановки любых двух столбцов, что и обеспечивает выполнение принципа Паули. Если бы две одночастичные функции совпадали, определитель обратился бы в нуль, что соответствует запрету на нахождение двух электронов в одном состоянии.

\subsection{Вычисление среднего значения гамильтониана}
\label{subsec:hf_energy_calculation}

Подставим функцию \eqref{eq:slater_det} в выражение для среднего значения гамильтониана. Одночастичная часть даёт сумму средних значений по занятым состояниям:
\begin{equation}
\label{eq:hf_one_electron}
\langle \hat{H}_1 \rangle = \sum_{a=1}^{Z} \int \dd[3]{\bm{r}} \psi_{\nu_a}^*(\bm{r}) \left( -\frac{\hslash^2}{2m} \nabla^2 - \frac{Z e^2}{r} \right) \psi_{\nu_a}(\bm{r}).
\end{equation}
Двухчастичная часть, описывающая межэлектронное взаимодействие, после учёта антисимметрии принимает вид
\begin{equation}
\label{eq:hf_two_electron}
\langle \hat{H}_2 \rangle = \frac{1}{2} \sum_{a,b=1}^{Z} \left( \langle a b | \hat{g} | a b \rangle - \langle a b | \hat{g} | b a \rangle \right),
\end{equation}
где $\hat{g} = e^2/\abs{\bm{r} - \bm{r}'}$, а матричные элементы определены как
\begin{equation}
\label{eq:matrix_elements}
\langle a b | \hat{g} | c d \rangle = \iint \dd[3]{\bm{r}} \dd[3]{\bm{r}'} \psi_{\nu_a}^*(\bm{r}) \psi_{\nu_b}^*(\bm{r}') \frac{e^2}{\abs{\bm{r} - \bm{r}'}} \psi_{\nu_c}(\bm{r}) \psi_{\nu_d}(\bm{r}').
\end{equation}
Первое слагаемое в \eqref{eq:hf_two_electron} описывает прямое кулоновское взаимодействие и совпадает с аналогичным членом в методе Хартри. Второе слагаемое --- обменное взаимодействие, которое возникает благодаря антисимметрии волновой функции. Обменный член отличен от нуля только для пар состояний с одинаковой проекцией спина, поскольку спиновые функции ортогональны:
\begin{equation}
\label{eq:spin_orthogonality}
\langle \chi_{\sigma_a} | \chi_{\sigma_b} \rangle = \delta_{\sigma_a, \sigma_b}.
\end{equation}
С учётом этого полная энергия в методе Хартри---Фока равна
\begin{equation}
\label{eq:hf_energy}
\begin{split}
E_{\text{HF}} &= \sum_{a=1}^{Z} \int \dd[3]{\bm{r}} \psi_{\nu_a}^*(\bm{r}) \left( -\frac{\hslash^2}{2m} \nabla^2 - \frac{Z e^2}{r} \right) \psi_{\nu_a}(\bm{r}) \\
&\quad + \frac{1}{2} \sum_{a,b=1}^{Z} \iint \dd[3]{\bm{r}} \dd[3]{\bm{r}'} \frac{e^2}{\abs{\bm{r} - \bm{r}'}} \abs{\psi_{\nu_a}(\bm{r})}^2 \abs{\psi_{\nu_b}(\bm{r}')}^2 \\
&\quad - \frac{1}{2} \sum_{a,b=1}^{Z} \delta_{\sigma_a, \sigma_b} \iint \dd[3]{\bm{r}} \dd[3]{\bm{r}'} \frac{e^2}{\abs{\bm{r} - \bm{r}'}} \psi_{\nu_a}^*(\bm{r}) \psi_{\nu_b}^*(\bm{r}') \psi_{\nu_b}(\bm{r}) \psi_{\nu_a}(\bm{r}').
\end{split}
\end{equation}

\subsection{Уравнения Хартри---Фока}
\label{subsec:hf_equations}

Варьируя функционал \eqref{eq:hf_energy} по $\psi_{\nu_a}^*(\bm{r})$ при условии ортонормированности одночастичных функций
\begin{equation}
\int \dd[3]{\bm{r}} \psi_{\nu_a}^*(\bm{r}) \psi_{\nu_b}(\bm{r}) = \delta_{ab},
\end{equation}
приходим к системе уравнений Хартри---Фока:
\begin{equation}
\label{eq:hartree_fock_eq}
\begin{split}
\left( -\frac{\hslash^2}{2m} \nabla^2 - \frac{Z e^2}{r} \right) \psi_{\nu_a}(\bm{r}) &+ \sum_{b=1}^{Z} \int \dd[3]{\bm{r}'} \frac{e^2 \abs{\psi_{\nu_b}(\bm{r}')}^2}{\abs{\bm{r} - \bm{r}'}} \psi_{\nu_a}(\bm{r}) \\
&- \sum_{b=1}^{Z} \delta_{\sigma_a, \sigma_b} \int \dd[3]{\bm{r}'} \frac{e^2}{\abs{\bm{r} - \bm{r}'}} \psi_{\nu_b}^*(\bm{r}') \psi_{\nu_a}(\bm{r}') \psi_{\nu_b}(\bm{r}) = \varepsilon_{\nu_a} \psi_{\nu_a}(\bm{r}).
\end{split}
\end{equation}
Первые два члена в левой части совпадают с уравнением Хартри. Третий член представляет собой нелокальный обменный потенциал. Из-за него уравнения Хартри---Фока существенно сложнее уравнений Хартри: они представляют собой систему интегро-дифференциальных уравнений, а не обычных дифференциальных.

\subsection{Обменное взаимодействие и фермиевская дырка}
\label{subsec:exchange_fermi_hole}

Обменный член удобно переписать в операторной форме:
\begin{equation}
\label{eq:exchange_operator}
-\sum_{b} \delta_{\sigma_a, \sigma_b} \int \dd[3]{\bm{r}'} \frac{e^2}{\abs{\bm{r} - \bm{r}'}} \psi_{\nu_b}^*(\bm{r}') \psi_{\nu_a}(\bm{r}') \psi_{\nu_b}(\bm{r}) = -\int \dd[3]{\bm{r}'} U(\bm{r}, \bm{r}') \psi_{\nu_a}(\bm{r}'),
\end{equation}
где введён нелокальный потенциал
\begin{equation}
\label{eq:nonlocal_potential}
U(\bm{r}, \bm{r}') = \frac{e^2}{\abs{\bm{r} - \bm{r}'}} \sum_{b} \delta_{\sigma_a, \sigma_b} \psi_{\nu_b}^*(\bm{r}') \psi_{\nu_b}(\bm{r}).
\end{equation}
Этот потенциал описывает обменное взаимодействие, которое не имеет классического аналога и обязано своим существованием принципу Паули. Можно показать, что обменное взаимодействие приводит к эффективному отталкиванию электронов с одинаковыми спинами. Это эквивалентно образованию так называемой фермиевской дырки --- области вокруг данного электрона, где вероятность найти другой электрон с той же проекцией спина мала.

Математически это выражается в том, что парная корреляционная функция для электронов с одинаковыми спинами обращается в нуль при совпадении координат:
\begin{equation}
\label{eq:fermi_hole}
\rho_2(\bm{r}, \bm{r}; \sigma, \sigma) = 0.
\end{equation}
Физически это означает, что два электрона с одинаковыми спинами не могут находиться в одной точке пространства. Именно это обстоятельство приводит к понижению энергии кулоновского отталкивания по сравнению с классическим случаем, поскольку электроны с параллельными спинами автоматически «раздвигаются» в пространстве.

Обменное взаимодействие можно интерпретировать и через матрицу плотности. Введём одночастичную матрицу плотности
\begin{equation}
\label{eq:density_matrix}
\rho(\bm{r}, \bm{r}') = \sum_{b} \psi_{\nu_b}^*(\bm{r}') \psi_{\nu_b}(\bm{r}).
\end{equation}
Тогда обменный потенциал записывается как
\begin{equation}
U(\bm{r}, \bm{r}') = \frac{e^2}{\abs{\bm{r} - \bm{r}'}} \rho(\bm{r}, \bm{r}').
\end{equation}
Для системы с полностью заполненными оболочками матрица плотности диагональна: $\rho(\bm{r}, \bm{r}') = \delta(\bm{r} - \bm{r}')$. В этом случае обменный потенциал принимает локальную форму, и уравнения Хартри---Фока упрощаются.

\subsection{Приближение Слэтера}
\label{subsec:slater_approximation}

На практике решение уравнений Хартри---Фока с нелокальным обменным потенциалом представляет собой сложную вычислительную задачу. Для упрощения часто используют приближение, предложенное Слэтером. Идея состоит в том, чтобы заменить нелокальный обменный потенциал некоторым локальным, зависящим только от электронной плотности.

В модели свободного электронного газа с плотностью $\rho$ матрица плотности имеет вид
\begin{equation}
\label{eq:free_electron_density_matrix}
\rho(\bm{r}, \bm{r}') = \frac{1}{2\pi^2 r^3} \left( \sin k_F r - k_F r \cos k_F r \right), \quad r = \abs{\bm{r} - \bm{r}'},
\end{equation}
где $k_F = (3\pi^2 \rho)^{1/3}$ --- волновой вектор Ферми. Усредняя обменный потенциал по всем электронам, получаем
\begin{equation}
\label{eq:slater_exchange}
V_{\text{обм}}(r) = -\frac{3}{2} \left( \frac{3}{\pi} \rho(r) \right)^{1/3} e^2.
\end{equation}
Это выражение известно как простая аппроксимация Слэтера. Оно часто используется в практических расчётах атомов, молекул и твёрдых тел, поскольку позволяет свести интегро-дифференциальные уравнения Хартри---Фока к обычным дифференциальным уравнениям.

\subsection{Теорема Купманса}
\label{subsec:koopmans_theorem}

Важным следствием уравнений Хартри---Фока является теорема Купманса. Умножим уравнение \eqref{eq:hartree_fock_eq} слева на $\psi_{\nu_a}^*(\bm{r})$ и проинтегрируем по $\bm{r}$. С учётом ортонормированности получаем
\begin{equation}
\label{eq:koopmans_derivation}
\varepsilon_{\nu_a} = \langle a | \hat{h} | a \rangle + \sum_{b} \left( \langle a b | \hat{g} | a b \rangle - \langle a b | \hat{g} | b a \rangle \right).
\end{equation}
Полная энергия при этом равна
\begin{equation}
\label{eq:total_energy_from_epsilon}
E_{\text{HF}} = \sum_{a} \varepsilon_{\nu_a} - \sum_{a < b} \left( \langle a b | \hat{g} | a b \rangle - \langle a b | \hat{g} | b a \rangle \right).
\end{equation}
Рассмотрим энергию, необходимую для удаления $a$-го электрона из атома. Если предположить, что волновые функции иона и атома одинаковы, то эта энергия равна взятой с обратным знаком одночастичной энергии $\varepsilon_{\nu_a}$. Это утверждение и составляет содержание теоремы Купманса:
\begin{equation}
\label{eq:koopmans}
I_a = -\varepsilon_{\nu_a},
\end{equation}
где $I_a$ --- потенциал ионизации $a$-й оболочки. Теорема Купманса позволяет интерпретировать собственные значения уравнений Хартри---Фока как потенциалы ионизации соответствующих состояний.

\section{Состояния электронов в самосогласованном центральном поле}
\label{sec:electron_states_central_field}

В приближении центрального поля одночастичный гамильтониан $\hat{h}_a$ коммутирует с операторами квадрата орбитального момента $\hat{\bm{l}}_a^2$ и его проекции $\hat{l}_{az}$, а также с операторами спина. Это означает, что состояние электрона характеризуется набором квантовых чисел $n, l, m, \sigma$. Энергия электрона $\varepsilon_{nl}$ зависит от $n$ и $l$, но не зависит от $m$ и $\sigma$. Кратность вырождения уровня $\varepsilon_{nl}$ по проекциям момента и спина равна $2(2l+1)$.

\subsection{Радиальное квантовое число и порядок уровней}
\label{subsec:radial_quantum_number}

Радиальное квантовое число $n_r$ равно числу нулей радиальной волновой функции. Главное квантовое число определяется как $n = n_r + l + 1$, как и в атоме водорода. Для заданного $l$ энергия возрастает с увеличением $n$. При фиксированном $n$ энергия тем ниже, чем меньше $l$, поскольку электроны с малым $l$ глубже проникают к ядру, где экранировка слабее, и сильнее притягиваются к ядру.

Это обстоятельство приводит к тому, что порядок уровней в многоэлектронном атоме отличается от водородного. В атоме водорода энергия зависит только от $n$, и уровни с одинаковым $n$, но разными $l$, вырождены. В самосогласованном поле это вырождение снимается. Например, уровень $3d$ лежит выше уровня $4s$, поскольку $4s$-электрон имеет большую вероятность находиться вблизи ядра и, следовательно, сильнее связан. Качественно это можно понять, рассмотрев эффективный потенциал: для электрона с малым $l$ центробежный барьер $\hslash^2 l(l+1)/(2mr^2)$ мал, поэтому он может подходить близко к ядру, где потенциал притяжения велик. Для электрона с большим $l$ центробежный барьер отталкивает его от ядра, и он движется в области, где экранировка сильнее, а притяжение слабее.

\subsection{Оболочки и подоболочки}
\label{subsec:shells_subs}

Совокупность состояний с данными $n$ и $l$ называют подоболочкой. Совокупность состояний с данным $n$ и всеми возможными $l$ (от $0$ до $n-1$) называют оболочкой. Число состояний в подоболочке с учётом спина равно $2(2l+1)$. В соответствии с принципом Паули на подоболочке может находиться не более $2(2l+1)$ электронов. Если подоболочка заполнена полностью, она называется замкнутой.

Оболочки принято обозначать буквами $K, L, M, N, \dots$, соответствующими значениям $n = 1, 2, 3, 4, \dots$. Внутри оболочки подоболочки обозначаются символами $s, p, d, f, \dots$, соответствующими значениям $l = 0, 1, 2, 3, \dots$. Например, подоболочка $3d$ соответствует $n = 3$, $l = 2$, а подоболочка $4s$ --- $n = 4$, $l = 0$.

\subsection{Правило Маделунга}
\label{subsec:madelung_rule}

Порядок заполнения оболочек в основном состоянии атома определяется последовательностью возрастания энергий $\varepsilon_{nl}$. В большинстве атомов он следующий:
\begin{equation}
\label{eq:filling_order}
1s \; 2s \; 2p \; 3s \; 3p \; 4s \; 3d \; 4p \; 5s \; 4d \; 5p \; 6s \; 4f \; 5d \; 6p \dots
\end{equation}
Энергии уровней, заключённых в близкие группы, могут меняться местами при переходе от одного атома к другому. Эмпирическое правило, сформулированное Маделунгом, позволяет предсказать порядок заполнения: заполнение идёт в порядке возрастания суммы $n+l$; при одинаковой сумме сначала заполняется уровень с меньшим $n$. Например, для $4s$ сумма $n+l = 4$, для $3d$ сумма $n+l = 5$, поэтому $4s$ заполняется раньше $3d$. Для $5s$ сумма $n+l = 5$, для $4d$ сумма $n+l = 6$, поэтому $5s$ заполняется раньше $4d$. Правило Маделунга имеет ряд исключений, особенно для тяжёлых атомов, но в целом хорошо описывает структуру периодической таблицы.

\subsection{Электронные конфигурации}
\label{subsec:electronic_configurations}

Распределение электронов по одночастичным состояниям называется электронной конфигурацией. Например, основное состояние гелия ($Z=2$) имеет конфигурацию $1s^2$, лития ($Z=3$) --- $1s^2 2s$, углерода ($Z=6$) --- $1s^2 2s^2 2p^2$. Запись $(nl)^k$ означает, что на подоболочке $nl$ находится $k$ электронов. Электроны, находящиеся на одной подоболочке, называются эквивалентными.

При заданной электронной конфигурации может существовать несколько различных состояний, отличающихся полным орбитальным моментом $L$ и полным спином $S$. Эти состояния называются атомными термами и будут подробно рассмотрены в следующей лекции. Здесь лишь отметим, что энергия этих состояний различна из-за остаточного кулоновского взаимодействия и обменных эффектов.

Химические свойства атома определяются в основном электронами внешних, незаполненных оболочек. Заполненные оболочки обладают сферически симметричным распределением заряда и слабо влияют на внешние электроны, экранируя заряд ядра. Именно поэтому атомы с заполненными внешними оболочками (инертные газы) химически инертны, а атомы с одним электроном на внешней оболочке (щелочные металлы) легко его теряют.

\section{Модель Томаса---Ферми}
\label{sec:thomas_fermi}

Для атомов с большим числом электронов ($Z \gg 1$) существует приближённый метод, основанный на квазиклассическом описании электронного газа. Этот метод был предложен Л. Томасом и Э. Ферми в 1927 году. Он позволяет найти усреднённую электронную плотность и потенциал, но не описывает оболочечную структуру.

\subsection{Основные предположения}
\label{subsec:tf_assumptions}

В модели Томаса---Ферми электроны рассматриваются как вырожденный газ свободных фермионов, находящихся в самосогласованном потенциале $\varphi(r)$. Предполагается, что потенциал достаточно медленно меняется с расстоянием, так что к рассматриваемому элементу объёма можно применять статистику Ферми---Дирака для свободных частиц. При абсолютном нуле температуры электроны заполняют все состояния с импульсами вплоть до импульса Ферми $p_F(r)$.

Условие того, что максимальная энергия электрона равна энергии Ферми $E_F$, имеет вид
\begin{equation}
\label{eq:fermi_energy}
\frac{p_F^2(r)}{2m} - e \varphi(r) = E_F.
\end{equation}
Для нейтрального атома начало отсчёта энергии выбирают так, что $E_F = 0$. Тогда
\begin{equation}
\label{eq:fermi_momentum}
p_F(r) = \sqrt{2 m e \varphi(r)}.
\end{equation}
Концентрация электронов связана с импульсом Ферми соотношением
\begin{equation}
\label{eq:electron_density_tf}
n(r) = \frac{p_F^3(r)}{3 \pi^2 \hslash^3} = \frac{(2 m e \varphi)^{3/2}}{3 \pi^2 \hslash^3}.
\end{equation}
Это соотношение получается из подсчёта числа квантовых состояний в фазовом пространстве: в объёме $\dd[3]{\bm{r}}$ с импульсами до $p_F$ число состояний равно $2 \cdot (4\pi p_F^3/3) \cdot \dd[3]{\bm{r}} / (2\pi \hslash)^3$, где множитель $2$ учитывает две проекции спина.

\subsection{Вывод уравнения Томаса---Ферми}
\label{subsec:tf_derivation}

Потенциал $\varphi(r)$ удовлетворяет уравнению Пуассона
\begin{equation}
\label{eq:poisson_tf}
\Delta \varphi = 4 \pi e n(r) - 4 \pi Z e \delta(\bm{r}).
\end{equation}
При $r \neq 0$ член с дельта-функцией отсутствует, и, подставляя \eqref{eq:electron_density_tf}, получаем
\begin{equation}
\label{eq:tf_equation_derivation}
\Delta \varphi = \frac{4 \pi e}{3 \pi^2 \hslash^3} (2 m e \varphi)^{3/2} = \frac{8 \sqrt{2} e}{3 \pi \hslash^3} (m e \varphi)^{3/2}.
\end{equation}
Это и есть уравнение Томаса---Ферми. Оно представляет собой нелинейное дифференциальное уравнение второго порядка для потенциала $\varphi(r)$.

Для решения этого уравнения удобно перейти к безразмерным переменным. Введём функцию $\chi(x)$ и переменную $x$ соотношениями
\begin{equation}
\label{eq:tf_dimensionless}
\varphi(r) = \frac{Z e}{r} \chi(x), \qquad r = b x, \qquad b = \frac{(3 \pi)^{2/3}}{2^{7/3}} \frac{\hslash^2}{m e^2} Z^{-1/3} \approx 0.885 a_{\text{B}} Z^{-1/3}.
\end{equation}
Подставляя эти выражения в \eqref{eq:tf_equation_derivation} и учитывая, что лапласиан в сферических координатах для функции, зависящей только от $r$, равен
\begin{equation}
\Delta \varphi = \frac{1}{r} \frac{d^2}{dr^2} (r \varphi),
\end{equation}
получаем универсальное уравнение Томаса---Ферми:
\begin{equation}
\label{eq:tf_universal}
\frac{d^2 \chi}{d x^2} = \frac{\chi^{3/2}}{\sqrt{x}}.
\end{equation}
Граничные условия следуют из условия нормировки и асимптотики потенциала:
\begin{equation}
\label{eq:tf_boundary}
\chi(0) = 1, \qquad \chi(\infty) = 0.
\end{equation}
Первое условие означает, что вблизи ядра потенциал стремится к кулоновскому потенциалу неэкранированного ядра. Второе условие соответствует нейтральному атому, у которого потенциал на больших расстояниях стремится к нулю.

\subsection{Численное решение и свойства функции $\chi(x)$}
\label{subsec:tf_numerical_solution}

Уравнение \eqref{eq:tf_universal} не зависит от $Z$, что является следствием универсальности модели. Функция $\chi(x)$ находится численным интегрированием. Она монотонно убывает и обращается в нуль только на бесконечности. Её асимптотика при больших $x$ имеет вид
\begin{equation}
\label{eq:tf_asymptotic}
\chi(x) \approx 144 x^{-3}, \quad x \gg 1.
\end{equation}
Эта асимптотика была найдена Зоммерфельдом. Она правильно передаёт поведение точного решения на больших расстояниях, хотя и не удовлетворяет граничному условию в нуле.

Зная $\chi(x)$, можно найти электронную плотность:
\begin{equation}
\label{eq:tf_density}
n(r) = \frac{Z^2}{4 \pi b^3} \left( \frac{\chi(x)}{x} \right)^{3/2}.
\end{equation}
Эта плотность обращается в нуль только на бесконечности, что отражает тот факт, что в модели Томаса---Ферми нейтральный атом не имеет резкой границы.

\subsection{Оценки характерных величин}
\label{subsec:tf_estimates}

Модель Томаса---Ферми позволяет получить простые оценки характерных величин. Определим средний радиус электронного облака $\langle r \rangle$ из условия, что внутри сферы этого радиуса заключена основная часть электронного заряда. Из вида потенциала \eqref{eq:tf_dimensionless} следует, что характерный масштаб длины в задаче равен $b \sim a_{\text{B}} Z^{-1/3}$. Поэтому
\begin{equation}
\label{eq:tf_radius}
\langle r \rangle \sim \frac{a_{\text{B}}}{Z^{1/3}}.
\end{equation}
Этот результат согласуется с оценкой, полученной ранее из качественных соображений. С ростом $Z$ размер атома уменьшается как $Z^{-1/3}$.

Средний импульс электрона оценим из соотношения неопределённостей $\langle p \rangle \sim \hslash / \langle r \rangle$:
\begin{equation}
\label{eq:tf_momentum}
\langle p \rangle \sim Z^{2/3} \frac{\hslash}{a_{\text{B}}}.
\end{equation}
Средняя кинетическая энергия электрона
\begin{equation}
\label{eq:tf_kinetic_energy}
\langle T \rangle \sim \frac{\langle p \rangle^2}{2m} \sim Z^{4/3} \, \text{Ry}.
\end{equation}
По теореме вириала для кулоновской системы средняя потенциальная энергия равна удвоенной средней кинетической энергии с обратным знаком:
\begin{equation}
\langle U \rangle = -2 \langle T \rangle.
\end{equation}
Поэтому средняя энергия одного электрона
\begin{equation}
\label{eq:tf_electron_energy}
E_e = \langle T \rangle + \langle U \rangle = -\langle T \rangle \sim -Z^{4/3} \, \text{Ry}.
\end{equation}
Полная энергия атома получается умножением на число электронов $Z$:
\begin{equation}
\label{eq:tf_atom_energy}
E_{\text{at}} \sim Z E_e \sim -Z^{7/3} \, \text{Ry}.
\end{equation}
Эта зависимость от $Z$ хорошо согласуется с более точными расчётами и экспериментальными данными для больших $Z$.

\subsection{Границы применимости}
\label{subsec:tf_limitations}

Модель Томаса---Ферми имеет ряд ограничений. Она не учитывает обменное взаимодействие и оболочечную структуру атома. Поэтому она применима для описания усреднённых характеристик атомов, но не для вычисления энергий ионизации и спектров. Наиболее надёжные результаты модель даёт для атомов с полностью заполненными оболочками и для усреднённых по оболочкам величин.

Кроме того, модель неприменима вблизи ядра, где потенциал меняется слишком быстро для квазиклассического приближения, и на больших расстояниях, где плотность электронов мала. Область применимости модели ограничена расстояниями
\begin{equation}
\frac{1}{Z} \ll \frac{r}{a_{\text{B}}} \ll 1.
\end{equation}
Тем не менее, модель Томаса---Ферми остаётся полезным инструментом для качественных оценок и для получения начальных приближений в более точных расчётах.

Учёт обменного взаимодействия в рамках статистической модели был выполнен Дираком. Соответствующее уравнение называется уравнением Томаса---Ферми---Дирака. Оно даёт более точные значения энергии, но качественная картина остаётся той же.

\end{document}
